\section{Reducci\'on de Feshbach--Schur}

Sea

\[
\mathcal H=\mathcal H_0\oplus\mathcal H_{\rm ext},
\qquad
H=
\begin{pmatrix}
H_{00}&V\\ V^*&H_{\rm ext}
\end{pmatrix}.
\]

Para \(z\) en el conjunto resolvente de \(H_{\rm ext}\), la compresi\'on del
resolvente es

\begin{equation}
P_0(z-H)^{-1}P_0
=\left[z-H_{00}-\Sigma_0(z)\right]^{-1},
\qquad
\Sigma_0(z)=V(z-H_{\rm ext})^{-1}V^*.
\label{rm:eq:rm-feshbach-schur}
\end{equation}

El operador \(\Sigma_0\) transporta a la hoja presente la memoria de las hojas
el\'iptica e hiperb\'olica. Su transformada temporal genera un n\'ucleo de
memoria

\[
S_{\rm eff}[\psi_0]
=\int\psi_0^\dagger(i\hbar\partial_t-H_{00})\psi_0\,\mathrm dt
-\iint\psi_0^\dagger(t)K_0(t,t')\psi_0(t')\,\mathrm dt\,\mathrm dt'.
\]

Una expansi\'on regular de baja frecuencia,

\[
\Sigma_0(\omega)
=\Sigma_0^{(0)}+\omega\Sigma_0^{(1)}
+\omega^2\Sigma_0^{(2)}+\cdots,
\]

produce una acci\'on local efectiva; una estructura espectral singular conserva
memoria no local. El presente recibe as\'i la respuesta del estado global sin
introducir un tiempo exterior.
